\documentclass[../../main.tex]{subfiles}

\begin{document}
\begin{CJK*}{UTF8}{gbsn}

\section*{Exercise 4}
Consider split conformal prediction 
with calibration data set $(X_1, Y_1), \cdots, (X_{n}, Y_{n}) 
\in \mathcal{X} \times \mathcal{Y}$ 
at given level $\alpha \in (0,1)$
and given score 
$s : \mathcal{X} \times \mathcal{Y} \to \mathbb{R}$.
We write $S_i = s(X_i, Y_i)$ for all $i$ for simplicity.
We look at the aggregating prediction sets by the majority vote 
$\mathcal{C}^{\text{mv}}(X_{n+1})$
of two conformal algorithms $\mathcal{C}_1$ and $\mathcal{C}_2$,
where $\mathcal{C}_1$ uses $s$ as score and 
$\mathcal{C}_2$ uses $-s$ as score.
Give an explicit formula of $\mathcal{C}^{\text{mv}}(X_{n+1})$
and its marginal coverage level, assuming exchangeability of the data set,
symmetry of $s$, and almost surely no 
ties in $S_i$'s.

\paragraph{Solution}
We first compute the expression of each conformal set.
Let us use the order statistics $S_{(1)} < \cdots < S_{(n)}$
of $S_1, \cdots, S_n$. 
Put 

\begin{equation*}
    k=\lceil(n+1)(1-\alpha)\rceil.
\end{equation*}
We discuss by cases of whether $k = n + 1$.

First of all, in the usual case, $ k < n + 1$.
The rule $(-v)_{(i)} = - v_{(n-i+1)}$ for any $v \in \mathbb{R}^n$
tells us that, by definition of $\mathcal{C}_1$ and $\mathcal{C}_2$:

\begin{equation*}
    \mathcal{C}_1(x)=\curly{y \in \mathcal{Y} \mid s(x,y) \leqslant S_{(k)} }
\end{equation*}
and 
\begin{equation*}
    \mathcal{C}_2(x)=\curly{y \in \mathcal{Y} \mid s(x,y) \geqslant S_{(n+1-k)}}.
\end{equation*}
Thus, by definition:

\begin{equation*}
    \mathcal{C}^{\text{mv}}(x)
    = \curly{y\in \mathcal{Y} \mid 
    \frac{1}{2} \sum_{i=1}^2 \chi_{y \in C_i(x)} > \frac{1}{2}} 
    =\mathcal{C}_1(x)\cap\mathcal{C}_2(x)
    =\curly{y\in \mathcal{Y} \mid S_{(n+1-k)}\leqslant s(x,y)\leqslant S_{(k)}}.
\end{equation*}
Next, for the trivial case $k = n + 1$, we have
$\mathcal{C}^{\text{mv}}(x)=\mathcal{C}_1(x)=\mathcal{C}_2(x)=\mathcal{Y}$.
Therefore:

\begin{equation*}
    \mathcal{C}^{\text{mv}}(X_{n+1})
    = 
    \begin{cases}
        \curly{y\in \mathcal{Y} \mid S_{(n+1-k)}\leqslant s(X_{n+1},y)\leqslant S_{(k)}}
        \tab[1.0cm] k=\lceil(n+1)(1-\alpha)\rceil < n + 1, \\
        \mathcal{Y} \tab[7.7cm] k= n + 1.
    \end{cases}
\end{equation*}
A simple counting of the probability of placing $S_{n+1}$ gives

\begin{equation*}
    P\curly{Y_{n+1} \in \mathcal{C}^{\text{mv}}(X_{n+1}) }
    =\frac{\max\{0,2k-n-1\}}{n+1}
    =\max\curly{0,1-\frac{2\lfloor\alpha(n+1)\rfloor}{n+1}}.
\end{equation*}
It just happens that this formula works for $k = n + 1$ case as well.
This completes the computation. ////

\end{CJK*}
\end{document}
